\documentclass[10pt]{article}
\linespread{1} 
\usepackage{extsizes}
\usepackage{amsmath}
\usepackage{amssymb}
\usepackage{amsthm}
\usepackage{ mathrsfs }
\usepackage{ mathtools}
\usepackage{enumerate}
\usepackage{bbm}
\usepackage{lipsum}
\usepackage{fancyhdr}
\usepackage{tikz-cd} 
\usepackage{adjustbox} 
\usepackage{ stmaryrd }
\numberwithin{equation}{section}
\theoremstyle{definition}
\usetikzlibrary{arrows}
\newcommand{\midarrow}{\tikz \draw[-triangle 90] (0,0) -- +(.1,0);}
\tikzset{commutative diagrams/.cd,
mysymbol/.style={start anchor=center,end anchor=center,draw=none}
}
\newcommand\comsymb[2][\circlearrowleft]{%
  \arrow[mysymbol]{#2}[description]{#1}}

\newtheorem{theorem}{Theorem} [section] 
\newtheorem{proposition}{Proposition}[section] 
\newtheorem{definition}{Definition}[section] 
\newtheorem{lemma}{Lemma}[section] 
\newtheorem{notation}{Notation}[section] 
\newtheorem{remark}{Remark}[section] 
\newtheorem{corollary}{Corollary} [section] 
\newtheorem{terminology}{Terminology}[section] 
\newtheorem{fact}{Fact}[section] 
\newtheorem{example}{Example}[section] 
\newtheorem{claim}{Claim}
\newenvironment{claimproof}[1]{\par\noindent\underline{Proof:}\space#1}{\hfill $\blacksquare$}

\DeclareMathOperator{\tmod}{mod}
\DeclareMathOperator{\triv}{triv}
\DeclareMathOperator{\ind}{ind}
\DeclareMathOperator{\straw}{s.t.}
\DeclareMathOperator{\ev}{ev}
\DeclareMathOperator{\pr}{pr}
\DeclareMathOperator{\Aut}{Aut}
\DeclareMathOperator{\Map}{Map}
\DeclareMathOperator{\Stab}{Stab}
\DeclareMathOperator{\Ind}{Ind}
\DeclareMathOperator{\GL}{GL}
\DeclareMathOperator{\SL}{SL}
\DeclareMathOperator{\SO}{SO}
\DeclareMathOperator{\Sp}{Sp}
\DeclareMathOperator{\esssup}{esssup}
\DeclareMathOperator{\abs}{abs}
\DeclareMathOperator{\rel}{rel}
\DeclareMathOperator{\diam}{diam}
\DeclareMathOperator{\rank}{rank}
\DeclareMathOperator{\Hom}{Hom}
\DeclareMathOperator{\Homk}{Hom_{k-alg}}
\DeclareMathOperator{\Ker}{Ker}
\DeclareMathOperator{\Image}{Im}
\DeclareMathOperator{\Dom}{Dom}
\DeclareMathOperator{\grad}{grad}
\DeclareMathOperator{\divergence}{div}
\DeclareMathOperator{\rk}{rk}
\DeclareMathOperator{\Span}{Span}
\DeclareMathOperator{\mSpec}{mSpec}
\DeclareMathOperator{\MaxSpec}{MaxSpec}
\DeclareMathOperator{\interior}{int}
\DeclareMathOperator{\supp}{supp}
\DeclareMathOperator{\id}{id}
\DeclareMathOperator{\sgn}{sgn}
\DeclareMathOperator{\Mat}{Mat}
\DeclareMathOperator{\PD}{PD}
\DeclareMathOperator{\ext}{ext}
\DeclareMathOperator{\vol}{vol}
\DeclareMathOperator{\inte}{int}
\DeclareMathOperator{\diverg}{div}
\DeclareMathOperator{\curl}{curl}
\DeclareMathOperator{\image}{im}
\DeclareMathOperator{\dR}{dR}
\DeclareMathOperator{\Ob}{Ob}
\DeclareMathOperator{\Mnfd}{Mnfd}
\DeclareMathOperator{\OpEmb}{OpEmb}
\DeclareMathOperator{\Spec}{Spec}
\DeclareMathOperator{\Spa}{Spa}
\DeclareMathOperator{\Lim}{Lim}
\DeclareMathOperator{\Alg}{\mathcal{A}}
\DeclareMathOperator{\glim}{glim}
\DeclareMathOperator{\con}{con}
\DeclareMathOperator{\Spv}{Spv}
\DeclareMathOperator{\clet}{c}
\DeclareMathOperator{\tat}{at}
\DeclareMathOperator{\Sheaf}{Sheaf}
\DeclareMathOperator{\Cont}{Cont}
%\newcommand*{\name}[\num_arguments][default values]{{\color{#1}\Large #2}}
\newcommand{\defeq}{\vcentcolon=}
\newcommand{\Ouv}{\mathscr{O}}
\newcommand{\norm}[1]{\Vert #1 \Vert}
\newcommand{\opNorm}[2]{\norm{#1}_{#2\to#2}}
\newcommand{\normL}[3]{\norm{#1}_{L^{#2}(#3)}}
\newcommand{\N}[0]{\mathbb{N}}
\newcommand{\m}[0]{\mathfrak{m}}
\newcommand{\R}[0]{\mathbb{R}}
\newcommand{\Z}[0]{\mathbb{Z}}
\newcommand{\C}[0]{\mathbb{C}}
\newcommand{\Q}[0]{\mathbb{Q}}
\newcommand{\bF}[0]{\mathbb{F}}
\newcommand{\U}[0]{\mathcal{U}}
\newcommand{\B}[0]{\mathcal{B}}
\newcommand{\G}[0]{\mathbb{G}}
\newcommand{\A}[0]{\mathbb{A}}
\newcommand{\T}[0]{\mathbb{T}}
\newcommand{\K}[0]{\mathcal{K}}
\newcommand{\Para}[0]{\mathbb{P}}
\newcommand{\M}[0]{\mathcal{M}}
\newcommand{\maniN}[0]{\mathcal{N}}
\newcommand{\cinf}[0]{\mathcal{C}^\infty}
\newcommand{\torus}[0]{\mathbb{T}}
\newcommand{\cR}[0]{\mathcal{R}}
\newcommand{\sep}[0]{\:|\:}
\newcommand{\pd}[2]{{\frac {\partial #1} {\partial #2}}}
\newcommand{\compinc}[0]{\subset \subset}
\newcommand{\sH}[0]{\mathscr{H}}
\newcommand{\Nfil}[0]{\mathcal{N}}
\newcommand{\g}[0]{\mathfrak{g}}
\newcommand{\p}[0]{\mathfrak{p}}
\newcommand{\QC}[0]{\mathfrak{Q}_C}
\newcommand{\ab}[1]{\vert #1 \vert}
\newcommand{\st}[0]{\:\straw\:}
\newcommand{\cGamma}[0]{c\Gamma}
\newcommand{\sC}[0]{\mathscr{C}}
\newcommand{\Grp}[0]{\mathrm{Grp}}
\newcommand{\op}[0]{\mathrm{op}}
\newcommand{\Func}[0]{\mathrm{Func}}
\newcommand{\sU}[0]{\mathscr{U}}
\newcommand{\cV}[0]{\mathcal{V}}
\newcommand{\sA}[0]{\mathscr{A}}
\newcommand{\sB}[0]{\mathscr{B}}
\newcommand{\sX}[0]{\mathscr{X}}
\newcommand{\sS}[0]{\mathscr{S}}
\newcommand{\sT}[0]{\mathscr{T}}
\newcommand{\cB}[0]{\mathcal{B}}
\newcommand{\cA}[0]{\mathcal{A}}
\newcommand{\cC}[0]{\mathcal{C}}
\newcommand{\cO}[0]{\mathcal{O}}
\newcommand{\cT}[0]{\mathcal{T}}
\newcommand{\cF}[0]{\mathcal{F}}
\newcommand{\cG}[0]{\mathcal{G}}
\newcommand{\nat}[0]{\mathrm{nat}}
\newcommand{\tand}[0]{\text{ and }}
\newcommand{\tor}[0]{\text{ or }}
\newcommand{\dcirc}[0]{{\circ\circ}}
\newcommand{\otherwise}[0]{\text{ otherwise }}
\newcommand{\Rat}[0]{\mathfrak{Rat}}
\newcommand{\polring}[1]{[X_1,\cdots,X_{#1}]}
\newcommand{\compring}[1]{\langle X_1,\cdots,X_{#1}\rangle}
\newcommand{\fib}[1]{%
  \mathbin{\mathop{\times}\limits_{#1}}%
}
\newcommand{\tens}[1]{%
  \mathbin{\mathop{\otimes}\displaylimits_{#1}}%
}
\newcommand{\comppol}[1]{\langle X_1,\cdots, X_{#1}\rangle}

\title{Linear Algebraic Groups}
\author{Based on Lectures given by\\\\Prof. Dr. Tasho Kaletha}
\date{WSSe 26/27, University of Bonn}


\begin{document}
\maketitle

%10/13
\section{Introduction}
% The course concerns with algebraic groups. Primary examples are 
% \begin{example}
%     $\GL_2(k)$ for any field $k$.
% \end{example}

% \begin{remark}
%     ALmost all finite simple groups arise as $G(k)/Z_G(k)$ for $k$ a finite field and $G$ a simple simply connected algebraic group.
% \end{remark}

\subsection{Definitions and Examples}

\begin{definition}
    For two objects $A_1,A_2\in\sC$ in a category, a product $(A,p_1:A\to A_1,A\to A_2)$ is a triplet consists of an object and two morphisms in $\sC$ with the following universal property. 
    For any object $B\in\sC$ together with morphisms $f_1:B\to A_1,f_2:B\to A_2$ in $\sC$, there is a unique morphism $B\to A$ making the following diagram commutes.
    \begin{center}
        \begin{tikzcd}
    & B \arrow[ldd, "f_1"'] \arrow[rdd, "f_2"] \arrow[d, "\exists!", dashed] &     \\
    & A \arrow[ld, "p_1"] \arrow[rd, "p_2"']                                 &     \\
A_1 &                                                                        & A_2
\end{tikzcd}
    \end{center}
    % is an object which is universal in terms of any object $B$ together with morphisms $f_1:B\to A_1,f_2:B\to A_2$.
    % That is, we have a unique morphism $B\to A$. We denote such $A$ to be $A_1\times A_2$. 
    By the universal property, such $A$ is unique up to isomorphisms. We denote such object as $A = A_1\times A_2$.
    \par A product of $n$-many objects are defined recursively by the above construction. If $n=0$, the product $1_\sC$ is called a terminal object and any $A\in\sC$, there is a unique morphism $1_A:A\to 1_\sC$. %(i.e. any object in $\sC$ admits a unique morphism to it). Such unique morphism is denote %We define a finite product for $n$-many object similarly. For $n=0$, such product is called the terminal object.
\end{definition}

% \begin{notation}
% As the terminal object $1_\sC$ admits a unique morphism into it from any object $A$, we denote such morphism as $1_A:A\to 1_\sC$.
% \end{notation}

\begin{remark}
    For $A\in\sC$, $A\times 1_\sC\to A$ given by the projection is an isomorphism with inverse $\id_A\times 1_A$.
\end{remark}

\begin{notation}
    For $A\in\sC$, $\Delta:A\to A\times A$ is induced from $A\overset{\id_A}{\underset{\id_A}{\rightrightarrows}} A$.
\end{notation}

\begin{definition}
    \label{definition_1_1}
    Let $\sC$ be a category admitting finite products. A group object in $\sC$ is a quintuple $(G,m,i,e)$ consists of an object $G\in\sC$ and three morphisms
    \begin{equation*}
        m:G\times G\to G, i:G\to G, e:1_\sC\to G,
    \end{equation*}
    satisfying the following axioms.
    \begin{enumerate}
        \item $m$ is associative. (i.e. $m\circ(m\times\id_G) = m\circ(\id_G\times m)$).
        \begin{center}
            \begin{tikzcd}
G\times G\times G \arrow[r, "\id_G\times m"] \arrow[d, "m\times\id_G"'] & G\times G \arrow[d, "m"] \\
G\times G \arrow[r, "m"']                                               & G                       
\end{tikzcd}
        \end{center}
        \item $e$ is a $2$-sided unit. (i.e. $m\circ(\id_G\times e) = \id_G = m\circ(e\times \id_G)$).
        \begin{center}
            \begin{tikzcd}
G \arrow[r, "e\times\id_G"] \arrow[rd, "\id_G"] \arrow[d, "\id_G\times e"'] & G\times G \arrow[d, "m"] \\
G\times G \arrow[r, "m"']                                                   & G                       
\end{tikzcd}
        \end{center}
        \item $i$ is a $2$-sided inverse. (i.e. $m\circ(i\times\id_G)\circ\Delta = e\circ 1_G = m\circ(\id_G\times i)\circ \Delta$).
        \begin{center}
            \begin{tikzcd}
                                                                      & G\times G \arrow[rd, "\id_G\times i"]  &                           \\
G \arrow[rd, "\Delta"'] \arrow[ru, "\Delta"] \arrow[r, "e\circ\id_G"] & G                                      & G\times G \arrow[l, "m"'] \\
                                                                      & G\times G \arrow[ru, "i\times \id_G"'] &                          
\end{tikzcd}
        \end{center}
    \end{enumerate}
\end{definition}

\begin{definition}
    \label{definition_1_2}
    A morphism of group objects is a morphism $f:G_1\to G_2$ satisfying 
    \begin{equation*}
        m_2\circ(f\times f)= f\circ m_1, i_2\circ f = f\circ i_1, e_2 = f\circ e_1.
    \end{equation*}
    The subclass of objects of groups together with these morphism forms a subcategory which we will denote by $\Grp\sC$.
\end{definition}

\begin{remark}
    \label{remark_1_3}
    For a group objects $G$ in $\sC$ and any object $A\in\Ob(\sC)$, $\sC(A,G)$ inhertis a group structure as folows. 
    \begin{enumerate}
        \item $e_A\defeq e\circ 1_A,A \to 1\to G$,
        \item $f_1\cdot f_2\defeq m\circ(f_1\times f_2)\circ\Delta_A$.
        \item $f^{-1}\defeq i\circ f$.
    \end{enumerate}
    Furthermore, the assingment $A\mapsto \sC(A,G)$ is a contravariant functor $\sC\to \Grp$ and $G\mapsto\sC(\cdot,G)$ is a contravariant functor $\Grp\sC\to \Func(\sC^\op,\Grp)$
    which is an embedding (i.e. fully-faithful) by Yoneda lemma.
    \par This is called the functor of points, we always have the abstract test group $G(1)$. In general, we will denote $G(A) \defeq \sC(A,G)$ for arbitrary $A\in\Ob(\sC)$.
\end{remark}

\begin{example}$\:$
    \begin{center}
\begin{tabular}{c|c|c}
Category $\mathscr{C}$ & $\mathrm{Grp}\mathscr{C}$ & Terminal Object \\ \hline
$\mathrm{Sets}$        & $\Grp$                    & $\{\ast\}$      \\ \hline
$\mathrm{Top}$         & $\mathrm{TopGrp}$         & $\{\ast\}$      \\ \hline
$\mathrm{Man}^\infty$  & $\mathrm{LieGrp}$         & $\{\ast\}$      \\ \hline
$\mathrm{ComplexMfd}$  & $\mathrm{ComplexLieGrp}$  & $\{\ast\}$      \\ \hline
$\mathrm{SmVar}_k$     & $\mathrm{AlgGrp}_k$       & $\Spec(k)$      \\ \hline
$\mathrm{SmVar}_S$     & $\mathrm{AlgGrp}_S$       & $S$            
\end{tabular}
    \end{center}
\end{example}

\begin{remark}
    For $\mathrm{Sets}$ and $G\in\Grp\mathrm{Sets}$, $G(\{\ast\})$ contains all the information about it. That is $G(\{\ast\}) = \Hom(\{\ast\},G)\cong G$.
\end{remark}

% \begin{example}
%     In $\sC = \mathrm{Sets}$, the terminal object is the singleton set and $\Grp\sC = \Grp$. $G(1)$ contains all the information about it.
% \end{example}
% \begin{example}
%     For $\sC = \mathrm{Top}$, the terminal object is again the singleton set and $\Grp\sC = \mathrm{Top}\Grp$. 
% \end{example}
% \begin{example}
%     For $\sC = \mathrm{Mani}^\infty$ the category of smooth manifolds, the terminal object is again the singleton set and $\Grp\sC = \mathrm{Lie}\Grp$. 
% \end{example}
% \begin{example}
%     For $\sC = \mathrm{CAMani}^\infty$ the category of complex analytic smooth manifolds, the terminal object is again the singleton set and $\Grp\sC = \mathrm{CALie}\Grp$ the category of complex analytic Lie groups.
% \end{example}
% \begin{example}
%     For $\sC = \mathrm{Sm}_k$ the smooth $k$-varieties for an algebraically closed field $k$, $\Grp\sC$ is the algebraic groups over $k$.
% \end{example}
% \begin{example}
%     For $\sC = \mathrm{Sch}/S$ the category of $S$-schemes, the terminal object is $S$ and$\Grp\sC$ is the category of $S$-group schemes.
% \end{example}

\begin{definition}
    \label{definition_1_4}
    For an algebraically closed field $k$, an affine algebraic $k$-group is a group object in the category of smooth (affine) $k$-varieties. 
    \par An (affine) $S$-group scheme is a group object in the category of $\mathrm{Sch}/S$.
\end{definition}

\begin{definition}
    \label{definition_1_5}
    Let $\sC$ be a category that admits finite co-products. A co-group object in $\sC$ is a group object in $\sC^\op$.
\end{definition}
%current Current
\begin{example}
    Let $\R$ be a commutative unital ring and $\sC$ the category of $R$-algebras which admits finite co-product. Then a co-group object in $\sC$ is an $R$-Hopf algebra. Explicitly, for $A$ an $R$-algebra, we have
    \begin{equation*}
        \delta:A\to A\tens{R}A, \varepsilon:A\to R, \eta:A\to A
    \end{equation*}
    which are called comultiplication, counit, and coinversion (or antipodal), respectively.
\end{example}
\begin{remark}
    Recall that we have an equivalence of categories,
    \begin{equation*}
        (R\mathrm{-Alg})^{\op}\stackrel{\Spec}{\to}\mathrm{AffSch}/R
    \end{equation*}
    inducing an equivalence between $R$-Hopf algebra and affine $R$-group schemes.
\end{remark}

\begin{remark}
    Let $G$ be an affine $k$-variety and $\cO_G$ be its ring of functions. Let $f\in\cO_G,f:G\to k$, then 
    \begin{equation*}
        f\circ m: G\times G\to k\in\cO_{G\times G}\cong \cO_G\tens{k}\cO_G.
    \end{equation*}
\end{remark}

\begin{example}
    Let $R$ be a commutative unital ring, $\underline{\mathbb{G}_m}$ be the functor such that 
    \begin{equation*}
        R\mathrm{-Alg}\ni S\stackrel{\underline{\mathbb{G}_m}}{\mapsto}(S^\times,\cdot)\in\Grp
    \end{equation*}
    is the functor of points of the $R$-group schemes with Hopf-algebra,
    \begin{equation*}
        A\defeq R[x,x^{-1}].
    \end{equation*}
    Structure morphisms are given by 
    \begin{equation*}
        A\ni x\stackrel{\delta}{\mapsto}x\otimes x\in A\otimes A, A\ni x\stackrel{\varepsilon}{\mapsto} 1_R\in R, x\stackrel{\eta}{\mapsto}x^{-1}\in A.
    \end{equation*}
\end{example}
\begin{example}
    Let $R$ be a commutative unital ring, $\underline{\mathbb{G}_a}$ be the functor such that 
    \begin{equation*}
        R\mathrm{-Alg}\ni S\stackrel{\underline{\mathbb{G}_m}}{\mapsto}(S,+)\in\Grp
    \end{equation*}
    is the functor of points of the affine $R$-group schemes with Hopf-algebra,
    \begin{equation*}
        A\defeq R[x].
    \end{equation*}
    Structure morphisms are given by 
    \begin{equation*}
        A\ni x\stackrel{\delta}{\mapsto}x\otimes 1_A+1_A\otimes x\in A\otimes A, A\ni x\stackrel{\varepsilon}{\mapsto} 0_R\in R, x\stackrel{\eta}{\mapsto}-x\in A.
    \end{equation*}
\end{example}

\begin{example}
    Let $R$ be a commutative unital ring, $\underline{\mathrm{GL}_N}$ be the functor such that 
    \begin{equation*}
        R\mathrm{-Alg}\ni S\stackrel{\underline{\mathrm{GL}_N}}{\mapsto}\GL_N(S)\in\Grp
    \end{equation*}
    is the functor of points of the affine $R$-group schemes with Hopf algebra,
    \begin{equation*}
        A = R[x_{11},\cdots,x_{NN},\det(x_{ij})^{-1}].
    \end{equation*}
    The structure morphisms are,
    \begin{align*}
        \delta(x_{ij}) &= \left(\sum_{k=1}^N x_{ik}x_{kj}\right),\\
        \varepsilon(x_{ij}) &= (\delta_{ij}),\\
        \eta(x_{ij}) &= (\det(x_{ij}))^{-1}\mathrm{adj}(x_{ij}).
    \end{align*}
\end{example}

\begin{example}
    $\underline{\SL_N}$ is the functor sending 
    \begin{equation*}
        R\mathrm{-Alg}\ni S\stackrel{\underline{\SL_N}}{\mapsto}\SL_N(S)\in\Grp.
    \end{equation*}
    Its Hopf algebra is 
    \begin{equation*}
        A = R[x_{ij}]_{ij}/(\det(x_{ij})-1).
    \end{equation*}
\end{example}
\begin{example}
    $\underline{\mu_N}$ is the functor sending 
    \begin{equation*}
        R\mathrm{-Alg}\ni S\stackrel{\underline{\mu_N}}{\mapsto}\{x\in S^\times\sep x^N=1\}\in\Grp.
    \end{equation*}
    Its Hopf algebra is 
    \begin{equation*}
        A = R[x]/(x^N-1)
    \end{equation*}
    and the structure morphisms are the same as in $\underline{\mathbb{G}_m}$.
\end{example}
\begin{example}
    $\underline{\alpha_p}$ for prime $p$ of the characteristic of $R$ is the functor sending 
    \begin{equation*}
        R\mathrm{-Alg}\ni S\stackrel{\underline{\alpha_N}}{\mapsto}\{x\in S^\times\sep x^p=0\}\in\Grp.
    \end{equation*}
    Its Hopf algebra is 
    \begin{equation*}
        A = R[x]/(x^p)
    \end{equation*}
    and the structure morphisms are the same as in $\underline{\mathbb{G}_a}$.
\end{example}
%OCT 15
\subsection{Basic Geometric Properties of Group Schemes}
\begin{definition}
    A scheme is separated if the diagonal map $\Delta:X\to X\times X$ is a closed immersion.
\end{definition}
\begin{remark}
    A topological space is Hausdorff if and only if $\Delta:X\to X\times X$ is closed. 
\end{remark}
\begin{lemma}
    \label{lemma_12}
    An $S$-group scheme $G$ is separated if and only if $e:S\to G$ is a closed immertion. This is automatic if $S = \Spec(k)$ or $G$ is affine.
\end{lemma}
\begin{proof}$\:$
    \begin{center}
        % https://tikzcd.yichuanshen.de/#N4Igdg9gJgpgziAXAbVABwnAlgFyxMJZABgBpiBdUkANwEMAbAVxiRAHEQBfU9TXfIRRkAjFVqMWbdgB0ZAMywAjYAGUunHn2x4CREeXH1mrRCFXdeIDDsH7SY6salnNVmwL0oATIaeTTc0ttTyFkAGY-CRNpYOt+XTDfR2iXDjiPRKJIlOdA2QVlNQ1ucRgoAHN4IlB5ACcIAFskXxAcCCRI1MDWLRB6pqQyNo7Ebz6B5sRh9qQDbrY5ABEYBhw6EGoGOiVVgAUEuzM6rAqACxw4ybnqWcQu7d2GA9svEBPzy-8Ys0a5AGMsHV-gAKORYKAAfXYpCwAEorg0pgAWW6jACs3zSvSs10QqJGSAAbFiepsQI99oc3h8LojBohMYTEAB2UmLGQrNYbCZI4lopBshZmMEyCHQuR4RrwAAEMDhAKB-xlImh5Mpz2pQnepzpXAoXCAA
\begin{tikzcd}
G \arrow[r] \arrow[d, "\Delta"']          & S \arrow[d, "e"] & S \arrow[r, "e"] \arrow[d, "e"']         & G \arrow[d, "\Delta"] \\
G\fib{S}G \arrow[r, "{m\circ(\id_G,i)}"'] & G                & G \arrow[r, "(\id_G\times (e\circ 1_G))"'] & G\fib{S}G            
\end{tikzcd}
    \end{center}
    Since closed immersion is stable under base change, the equivalence is done.
\end{proof}
\begin{remark}
    A topological gropu$G$ is Hausdorff if and only if $e\in G$ is a closed point (i.e. $\{e\}$ is closed).
\end{remark}
Recall that $X$ scheme is connected or irreducible if $\ab{X}$ has that property, where $\ab{X}$ is the underlying space. In general, irreducibility implies the connectedness. 
\begin{lemma}
    \label{lemma_13}
    For a $S$-group scheme, a connected component is an irreducible component. In particular, these two notions are equivalent.
\end{lemma}
\begin{proof}
    Since an irreducible component is not fully contained in the union of two others, there is a point contained only on an irreducible component.
    By translation every point lies in a unique irreducible component, thus the irreducible components are disjoint.
\end{proof}
\begin{remark}
    If $G$ is smooth affine, irreducibility, connectedness and $\cO_G$ being an integral domain are all equivalent.
\end{remark}
\par\textbf{Identity component}
Let $X$ be a scheme and $Z\subseteq\ab{X}$ be a closed subset, let $Z\hookrightarrow X$ be a canonical inclusion as topological spaces. Then $(Z,i^*\cO_X)$ is a locally ringed space but may not be a scheme.
To make this a scheme we need $(Z,i^*(\cO_X/\mathcal{I}_Z))$ where 
\begin{equation*}
    \mathcal{I}_Z(U) = \{f\in\cO_X(U)\sep \forall x\in U\cap Z, f(x) = 0 \text{ in $\cO_{X,x}/\m_x$}\}.
\end{equation*}
This is called the reduced scheme structure on $Z$, however, if $Z$ is generalizing (closed under generalization) then $(Z,i^*\cO_X)$ is a scheme.

\begin{lemma}
    \label{lemma_14}
    A connected component is generalizing.
\end{lemma}
\begin{proof}
    Since $Z$ and $\overline{\{y\}}$ are both connected and meet at $x$, $Z\cup\overline{\{y\}}$ is also connected hence $Z = Z\cup\overline{\{y\}}$ thus $y\in Z$.
\end{proof}
\begin{definition}
    \label{definition_15}
    For a $k$-group scheme $G$, the identity component $G^0$ is defined as the scheme $(\ab{G}^0,\cO_G|_{\ab{G}^0})$ which is a group scheme.
\end{definition}
Recall if $(R,\m)$ is a local ring, then the dimension of $R$ is the maximal chain of prime ideals which is called the Krull-dimension. Also $\m/\m^2$ is a $k$-vector space where $k=R/\m$, hence we can consider the dimension of $\m/\m^2$ as a $k$-vectorspace. $R$ is said to be regular if these two dimensions coincide.
\begin{definition}
    A scheme $X$ is regular if all local rings $\cO_{X,x}$ are regular.
\end{definition}
\begin{definition}
    A $k$-scheme $X$ of finite type is smooth if $X\fib{k}\overline{k}$ is regular.
\end{definition}
\begin{definition}
    For $S$-scheme $X$, geometric fibers are the base change $X_{\overline{S}}$
    \begin{center}
        \begin{tikzcd}
X_{\overline{S}} \arrow[d] \arrow[r] & X \arrow[d] \\
\Spec(k) \arrow[r]                   & S          
\end{tikzcd}
    \end{center}
    where $k$ is algebraically closed.
\end{definition}
\begin{definition}
    A $S$-scheme $X$ is smooth if it is flat, locally finite presentation and its geometric fibers are regular. 
\end{definition}
\begin{lemma}
    \label{lemma_16}
    For the notion of smoothness for group scheme, the geometric fibers being regular can be replaced by the following criteria. Every point $\overline{e}$ of a geometric base change $X_{\overline{S}}$ over $e$, $X_{\overline{S}}$ is regular at $\overline{e}$.
\end{lemma}
\begin{proof}
    Clearly if $G$ is smooth, the above condition is met. Conversely, for any $g\in G(k)$, where $k$ is algebraically closed, we have
    \begin{equation*}
        L_g:G\fib{S}\Spec(k)\stackrel{\sim}{\to} G\fib{S}\Spec(k).
    \end{equation*}
    Therefore, $\cO_{G,g}$ is regular. The smooth locus $U\subseteq G$ is open, non-empty. Take $V= G\backslash U$. If $V$ is not empty then some affine open has a $k$-point for some algebraically closed field $k$ by Nullstellensatz hence a contradiction.
\end{proof}
\begin{definition}
    A scheme $X$ is reduced if all its stalks are reduced as rings. Equivalently all affine opens are spectrum of reduced rings.
\end{definition}
\begin{remark}
    Even though $X$ is reduced, for a morphism of schemes $X\to Y$, the fibers might not be reduced. The following is a counter example,
    \begin{equation*}
        (\Z[x]/(x^p-1))\otimes \bF_p.
    \end{equation*}
\end{remark}
\begin{remark}
    $X/k$ is reduced but $X\fib{k}\overline{k}$ needs not be reduced. We can construct a counter example as follows.
    \begin{equation*}
        k = \bF_p(s), R = k[t]/(t^p-s)
    \end{equation*}
    is reduced however,
    \begin{equation*}
        R\tens{k}k(\sqrt[p]{s})
    \end{equation*}
    is not.
\end{remark}
For a scheme $X$ over $k$, geometrically reduced implies reduced. If $k$ is perfect then these two notions coincide.
\begin{lemma}
    \label{lemma_17}
    A $k$-group scheme of locally finite type is smooth if and only if it is geometrically reduced.
\end{lemma}
\begin{proof}
    By the definition of smoothness, we assume $k$ to be algebraically closed without loss of generality.
    Recall that regular local ring is reduced. Conversely, a smooth scheme $X$ is reduced which is locally of finite type, then it has non-empty smooth locus which is open and dense.
    If $X=G$ a group scheme then using translation, all points are smooth.
\end{proof}
\begin{proposition}
    \label{proposition_18}
    For a $k$-group scheme $G$ locally of finite type where $k$ is algebraically closed, $G$ satisfies the following properties.
    \begin{enumerate}[1. ]
        \item $G$ is equidimensional and $\dim(G) \leq \dim T_e(G)$.
        \item The following statements are equivalent for $G$
        \begin{enumerate}[i. ]
            \item $G$ is smooth.
            \item $G^0$ is smooth.
            \item $\cO_{G,e}$ is regular.
            \item $\dim(G) = \dim T_e(G)$.
            \item $G$ is reduced.
            \item For any $k$-algebra $R$ and $I\subseteq R$ an ideal such that $I^2 = 0$, $G(R)\twoheadrightarrow G(R/I)$.
        \end{enumerate}
    \end{enumerate}
\end{proposition}
\begin{proof}
    $G$ is equi-dimensional by translation. By Nakayama's lemma $\dim T_e(G)\geq \dim(\cO_{G,e})$. 
    \par For i $\iff$ ii $\iff$ iii are due to translation. iii $\iff$ iv is by definition. i$\iff$ v is by Lemma \ref{lemma_14}. i$\iff$vi by lifting criterion.
\end{proof}
Let $R$ be a finite $k$-algebra and $X=\Spec(R)$ where $k$ is algebraically closed. Since $x\in X_{\mathrm{cl}}$ corersponds to a maximal ideal in $R$, the composition,
\begin{equation*}
    k\to R\to R/\m_x
\end{equation*}
is an isomorphism. We write $f(x) = \ev_x(f)\in k$. 

\begin{definition}
    A derivation $d:R\to k$ is a $k$-linear map satisfying Leibniz rule. 
\end{definition}
In the above setting, we have 
\begin{equation*}
    T_xX = \mathrm{Der}_k(R,k)
\end{equation*}
and restriction to $\m_x$ gives an isomorphism of $\mathrm{Der}_k(R,k)$ and $\Hom_k(\m_x/\m_x^2,k)$.
\par A vector field is $d:x\mapsto d_x$ such tha for each $f\in R$, $d(f)\in R$. We have a morphism,
\begin{equation*}
    \mathrm{Der}_k(R,R)\stackrel{d\mapsto \ev_x\circ d}{\to}\mathrm{Der}_k(R,k).
\end{equation*}
\begin{proposition}[Cartier]
    Let $k$ be a field of characteristic $0$. Any (affine) k-group scheme of finite type is automatically smooth.
\end{proposition}
\begin{proof}
    As usual, we assume $k$ is algebraically closed. Let $I\subseteq\cO_G$ be nilradical, then $I\subseteq\m_{G,e}\subseteq\cO_G$.
    \begin{claim}
    $I\subseteq\m_{G,e}^2$
    \end{claim}
    \begin{claimproof}
        Let $d_e\in T_eG=\mathrm{Der}_k(\cO_G,k)=\Hom_k(\m_{G,e}/\m_{G,e},k)$. Hence enough to show that $d_e(I) = 0$. Take $d\in\mathrm{Der}_k(\cO_{G,e},\cO_{G,e})$ that lifts $d_e$, using Leibniz rule we see $d_e(I)=0$. Indeed for $f\in I$, if $N$ is the smallest natural number such that $f^N = 0$, then 
        \begin{equation*}
            d_e(f^N) = Nf^{N-1}d(f),
        \end{equation*}
        Note that $N$ is not $0$ as $k$ is of characteristic $0$ and $f^{N-1}$ is not $0$. Hence $d(f)$ is not a unit. Hence $d(f)\in\m_{G,e}$. Recall we have $d_e(f) = \ev_e(d(f))=0$. 
    \end{claimproof}
    Let $H=G_{\mathrm{red}}$ which is a reduced $k$-group scheme hence smooth by Lemma % cite the right lemma
    We have $\cO_H=\cO_G/I$, hence $\m_{H,e} = \m_{G,e}/I$. With the claim we have $\m_{H,e}^2 = \m_{G,e}^2/I$. Therefore $\m_{H,e}/\m_{H,e}^2\cong \m_{G,e}/\m_{G,e}^2$. Now $\dim G\leq \dim T_e(G)=\dim T_e(H) = \dim(H)=\dim G$ by Proposition \ref{proposition_18}. Hence $\dim G = \dim T_e(G)$ hence $G$ is smooth.
\end{proof}

%10/20

\subsection{Kernels and Exact Sequences}

\begin{definition}
    \label{definition_21}
    The kernel of a homomorphism of algebraic group schemes $\varphi:G\to H$ is the fiber product 
    \begin{center}
        \begin{tikzcd}
\ker\varphi\defeq G\fib{H}\{\ast\} \arrow[d] \arrow[r] & \{\ast\} \arrow[d] \\
G \arrow[r, "\varphi"']                               & H                 
\end{tikzcd}
    \end{center}
\end{definition}

\begin{remark}
    Obviously $\ker\varphi(R)=\ker(\varphi(R):G(R)\to H(R))\trianglelefteq G(R)$ and hence $\ker\varphi\trianglelefteq G$.
\end{remark}

\begin{remark}
    If $G,H$ are both affine then so is $\ker\varphi$ with the underlying Hopf algebra 
    \begin{equation*}
        \cO_{\ker\varphi} = \cO(G)\tens{\cO(H)}k\cong \cO(G)/I_H\cO(G),
    \end{equation*}
    where $I_H \defeq \ker(\cO(H)\to k)$ which is called the augumentation ideal, (i.e., the kernel of morphism $\cO(H)\ni f\mapsto f(e)\in k$).
\end{remark}

\begin{example}
    $\mu_n = \ker(\G_m\stackrel{(\cdot)^n}{\to}\G_m)$ and $\alpha_p = \ker(\G_A\stackrel{(\cdot)^p}{\to}\G_A)$ where the latter group scheme is over a field of characteristic $p$.
\end{example}

\begin{definition}
    A morphism of schemes is faithfully-flat if it is flat and surjective as a morphism between underlying topological spaces.
\end{definition}

\begin{definition}
    \label{definition_23}
    The sequence 
    \begin{equation*}
        e\to N\stackrel{i}{\to}G\stackrel{q}{\to}Q
    \end{equation*}
    is exact if $i$ induces an isomorphism between $N$ and $\ker q$.
    \par The sequence 
    \begin{equation*}
        e\to N\stackrel{i}{\to}G\stackrel{q}{\to}Q\to e
    \end{equation*}
    is exact if in addition $q$ is faithfully-flat.
\end{definition}

\begin{proposition}
    \label{proposition_24}
    Given an exacr sequence,
    \begin{equation*}
        e\to N\to G\to Q\to e,
    \end{equation*}
    \begin{enumerate}
        \item  $N,Q$ are smooth $\Rightarrow$ $G$ is smooth.
        \item $G$ is smooth $\Rightarrow$ $Q$ is smooth.
    \end{enumerate}
\end{proposition}


\begin{proof}
    \par\textbf{1.} If $q$ is flat and fibers of $q$ are the translates of $N$ hence smooth. Note that $q$ is smooth if $N$ is smooth and hence $G$ is smooth as $Q$ is.
    \par\textbf{2.} If $q$ is faithfully flat then $\cO_Q\hookrightarrow q_*\cO_G$ is injective. As $G$ is smooth, it is geometrically reduced, so is $\cO_Q$ hence $Q$ is smooth.
\end{proof}

\begin{example}
    An example of $G$ being smooth but $N$ fails to be is $\G_m$ and $\mu_p$ where the involving field is of characteristic $p$.
\end{example}
\begin{proposition}
    \label{proposition_25}
    Let $\varphi:G\to H$ be a surjective morphism of group varieties. The following statements about $\varphi$ are equivalent.
    \begin{enumerate}[1. ]
        \item $d_e\varphi:T_eG\to T_eH$ is surjective.
        \item $\ker\varphi$ is smooth.
        \item $\varphi$ is smooth.
    \end{enumerate}
\end{proposition}
\begin{proof}
    WIthout loss of generality, we assume $k$ is algebraically closd. Write $N=\ker \varphi$. Recall the equality agian that the fibers of $\varphi$ are the translates of $N$ hence all dimensions of fibers are equal to $\dim(N)$. In particular, we obtain,
    \begin{equation*}
        \dim(G) = \dim(H)+\dim(N).
    \end{equation*}
    Recall by Proposition \ref{proposition_18}
    \begin{equation*}
        \dim T_eG \geq \dim G.
    \end{equation*}
    The equality holds if and only if $G$ is smooth. Consider 
    \begin{center}
        \begin{tikzcd}
0 \arrow[r] & T_eN \arrow[d] \arrow[r]              & T_eG \arrow[d] \arrow[r, "d_e\varphi"] & T_eH \arrow[d]              \\
0 \arrow[r] & N(k(\varepsilon)) \arrow[d] \arrow[r] & G(k(\varepsilon)) \arrow[d] \arrow[r]  & H(k(\varepsilon)) \arrow[d] \\
0 \arrow[r] & N(k) \arrow[r]                        & G(k) \arrow[r]                         & H(k)                       
\end{tikzcd}
    \end{center} 
    We have $\dim T_eG \leq \dim T_eN+\dim T_eH$ and the equality holds if and only if $\varphi$ is surjective. Using once again Proposition \ref{proposition_18}, we see 
    \begin{equation*}
        \dim G\leq \dim T_eN+\dim H,
    \end{equation*}
    hence \textbf{i}$\iff$\textbf{ii}. 
    \par \textbf{ii}$\iff$\textbf{iii}. $\varphi$ smooth means fibers are smooth hence $\Leftarrow$. Conversely, $d_e\varphi$ being surjective means that $d_g\varphi$ is surjective for all $g\in G(k)$ by the homogeneity, and as $G,H$ are smooth so is $\varphi$.
\end{proof}

\subsection{Group Actions}
\begin{definition}
    \label{definition_26}
    An action of a $k$-group scheme $G$ on a $k$-scheme $X$ is a regular map $\mu:G\times X\to X$ that the following diagrams are commutative.
    \begin{center}
        \begin{tikzcd}
G\times G\times X \arrow[d, "m\times\id_X"'] \arrow[r, "\id_G\times \mu"] & G\times X \arrow[d, "\mu"] & \{\ast\}\times X \arrow[r, "e\times\id_X"] \arrow[rd, "\cong"'] & G\times X \arrow[d, "\mu"] \\
G\times X \arrow[r, "\mu"']                                               & X                          &                                                                 & X                         
\end{tikzcd}
    \end{center}
    We will write $g\cdot x, gx, g\times x$ for $\mu(g,x)$. Giving group scheme actions is equivalent to giving actions $G(R)\circlearrowright X(R)$.
\end{definition}
\begin{lemma}
    \label{lemma_27}
    Consider $G\stackrel{\mu}{\circlearrowright}X$ which is faithfully flat. Then we have the following statements about $\mu$.
    \begin{enumerate}
        \item $G$ is smooth then $\mu$ is smooth.
        \item $G$ is finite then $\mu$ is finite.
    \end{enumerate}
\end{lemma}
\begin{proof}
    If $(g,x)\mapsto (g,g\cdot x)$ is an isomorphism then by the diagram 
    \begin{center}
        \begin{tikzcd}
G\times X \arrow[d, "\mu"'] \arrow[r, "{(g,x)\mapsto (g,g\cdot x)}"] & G\times X \arrow[d, "p_X"] \\
X \arrow[r, "\id_X"']                                                & X                         
\end{tikzcd}
    \end{center}
    any properties of $\mu$ hold for $p_2$ and therefore for $G$ as well.
\end{proof}
\begin{definition}
    A morphism of $k$-schemes is $G$-equivariant if it commutes with $G$-actions.
\end{definition}
\begin{proposition}
    \label{proposition_28}
    $X\stackrel{f}{\to}Y$ be $G$-equivariant between non-empty $k$-schemes.
    \begin{enumerate}
        \item $Y$ is reduced and $G(\overline{k})\circlearrowright Y(\overline{k})$ is transitive then $f$ is faithfully flat.
        \item If $G(\overline{k})\circlearrowright X(\overline{k})$ is transitive then $f(\ab{X})\subseteq\ab{Y}$ locally closed and hence $f(X)_{\mathrm{red}}\defeq f(\ab{X})$ with its reduced scheme structure.
        \item If $G(\overline{k})\circlearrowright X(\overline{k})$ is transitive and $X$ is reduced then $f$ factors as 
        \begin{equation*}
            X\stackrel{\text{faithfully flat}}{\to}f(X)_{\mathrm{red}}\hookrightarrow Y,
        \end{equation*}
        and $f(X)_{\mathrm{red}}$ is $G$-stable.
    \end{enumerate}
\end{proposition}
\begin{proof}
    See \cite{Milne}.
\end{proof}

\begin{definition}
    \label{definition_29}
    Let $G$ be a group scheme and $X$ be a scheme where $G$ acts on. For $x\in X(k)$, we have the orbit map is 
    \begin{equation*}
        G\ni g\stackrel{\mu_x}{\mapsto} gx\in X.
    \end{equation*}
    $\mu_x$ is the restriction of $\mu$ to $G\times\{x\}$. If $G(\overline{k})$ acts transitively on $X(\overline{k})$ then $\mu_x$ is surjective.
    \par The orbit $G\cdot x$ is the image of $\ab{\mu_x}$. 
\end{definition}

\begin{remark}
    Due to Proposition \ref{proposition_28}.2, $\ab{\mu_x}$ is locally closed together with considering the trivial action of $G$ on itself.
\end{remark}

\begin{example}
    \label{example_30}
    \begin{equation*}
        \SL_2\times\A \ni\begin{pmatrix}
            a & b \\ c & d 
        \end{pmatrix}
        \begin{pmatrix}
            x \\ y 
        \end{pmatrix}
        = \begin{pmatrix}
            ax + by \\ cx+dy
        \end{pmatrix}
    \end{equation*}
    Then we have two orbits,
    \begin{equation*}
        \{0\},\left\{\begin{pmatrix}
            x \\ y
        \end{pmatrix}\not=0\right\}.
    \end{equation*}
    The first one is clearly closed however, the latter has the closure the whole hence not closed even it is locally closed.
\end{example}
\begin{proposition}
    \label{proposition_31}
    Let $G$ be a $k$-group variety acting on $k$-variety $X$ where $k$ is algebraically closed. Then every orbit of minimal dimension is closed. In particular, there will always be a closed orbit.
\end{proposition}
\begin{proof}
    Let $O\defeq G\cdot x$ be an orbit. Then there is $U\subseteq \overline{O}$ a dense open subset contained in $O$ by Chevalley's theorem. Observe that 
    \begin{equation*}
        O = \bigcup_{g\in G(k)}gU \subseteq \overline{O},
    \end{equation*}
    hence open. Therefore $\overline{O}\backslash O$ is closed and is of dimension less than $\dim O$. If $O$ is closed then there is nothing to show.
    Otherwise $\overline{O}\backslash O\not=\emptyset$ is a union of orbit of dimension less than $\dim O$ hence $O$ is closed if it is of minimal dimension.
\end{proof}

\begin{theorem}[Homomorphism Theorem]
    \label{theorem_32}
    Let $G\stackrel{\varphi}{\to}H$ be a homomorphism of group schemes. if $G$ is smooth then $\varphi$ factors as,
    \begin{equation*}
        G\stackrel{\text{faithfully flat}}{\to} I\stackrel{\text{closed immersion}}{\to}H.
    \end{equation*}
\end{theorem}

\begin{thebibliography}{9}
\bibitem{Borel} Linear algebraic groups.
\bibitem{Milne} Linear algebraic groups.
\end{thebibliography}
\end{document}